%% Shared manuscript BODY -- this is the file you edit.
%% Compiled by both submission.tex (single-column) and editing.tex (two-column).
%% Mapped to notes/paper.nature-biotech-cgt-outline.md.

%% ===== Title (aim <= ~15 words) =====
\title[TorchCell CGT for multimodal phenotype prediction]{Accelerating metabolic engineering through multimodal phenotype prediction with TorchCell and the Cell Graph Transformer}

%% ===== Authors / affiliations (FILL IN) =====
\author*[1]{\fnm{Michael} \sur{Volk}}\email{mjvolk3@illinois.edu}
\author[1]{\fnm{First} \sur{Coauthor}}\email{}
\author*[1]{\fnm{Huimin} \sur{Zhao}}\email{zhao5@illinois.edu}

\affil[1]{\orgdiv{Department of Chemical and Biomolecular Engineering}, \orgname{University of Illinois Urbana-Champaign}, \orgaddress{\city{Urbana}, \state{IL}, \country{USA}}}

%% ===== Abstract (<= 150 words; shortened from SIMB 2026) =====
%% Point estimates only (no +/-) to sidestep the unresolved error-bar provenance
%% (see outline Gaps); restore uncertainties once reconciled.
\abstract{Deep learning has advanced protein structure prediction, yet metabolic
engineering lags because phenotype data are scattered across heterogeneous studies
lacking a shared ontology. We present
TorchCell, an open-source framework that annotates literature and public datasets
with a shared ontology and ingests them into a Neo4j knowledge graph, then builds
deep-learning-ready datasets in which each strain is a perturbation operator over
a shared wildtype reference cell combining genome sequence, gene networks,
metabolism, and environment. We develop the Cell Graph Transformer (CGT), a
virtual-cell architecture in which biological interaction graphs constrain
multi-head attention and an equivariant perturbation operator feeds multitask
decoders. In \textit{Saccharomyces cerevisiae}, CGT predicts trigenic
gene-gene interactions (Pearson $r=0.454$), outperforming DANGO, DCell, and Yeast9
flux balance analysis, and generalizes to knockout expression ($r=0.543$) and
morphology ($r=0.619$). CGT recommends gene deletions for beta-carotene and
betaxanthin production, pairing with autonomous biofoundries for AI-guided strain
engineering.}

\keywords{metabolic engineering, foundation models, gene interactions, virtual cell, knowledge graph, strain design}

\maketitle

%% ===== Introduction (~500-650 words) =====
\section{Introduction}\label{sec:intro}
\textit{[FILLER -- replace with Intro prose.]} Deep learning transformed protein
structure and function prediction, yet systems biology and metabolic engineering
lag because phenotype data are scattered across small, heterogeneous studies with
no shared schema~\cite{volk2023meteng}. Prior approaches fall short: mechanistic
flux balance analysis misses epistasis, and recent perturbation transformers do
not consistently beat linear baselines~\cite{ahlmanneltze2025}. Our own classical
baselines predict knockout fitness but fail on gene interactions (Suppl.\ Fig.\
S1), motivating a deep, biologically grounded model. Here we present TorchCell and
the Cell Graph Transformer (CGT).

%% ===== Results (~1,700-2,000 words) =====
\section{Results}\label{sec:results}

\bmhead{A shared ontology and knowledge graph unify metabolic-engineering data}
\textit{[FILLER -- R1, see Fig.~\ref{fig:overview}.]} TorchCell annotates
literature and public datasets with a shared experiment ontology and ingests them
into a Neo4j knowledge graph; dataloaders emit deep-learning-ready datasets in
which each strain is a perturbation operator over a shared wildtype reference cell.

\bmhead{The Cell Graph Transformer: a graph-constrained virtual cell}
\textit{[FILLER -- R2.]} Biological interaction graphs constrain multi-head
attention through a loss aligning attention to network priors; an equivariant
perturbation operator maps the reference embedding to the perturbed strain state.

\bmhead{CGT predicts trigenic gene-gene interactions at state of the art}
\textit{[FILLER -- R3.]} Classical ML fails on this task (Suppl.\ Fig.\ S1); CGT
reaches Pearson $r = 0.454 \pm 0.006$, outperforming DANGO, DCell, and Yeast9 FBA.

\bmhead{One embedding, many phenotypes: expression and morphology}
\textit{[FILLER -- R4.]} The same architecture predicts knockout expression
($r = 0.543$) and single-knockout morphology ($r = 0.619$).

\bmhead{CGT-guided strain design with an autonomous foundry}
\textit{[FILLER -- R5.]} We recommend gene deletions for beta-carotene and
betaxanthin production and pair with the UIUC iBioFoundry for iterative design.

%% Full-width figure (figure* spans both columns in two-column; full width in single).
\begin{figure*}[t]
\centering
\fbox{\parbox[c][4.5cm][c]{0.95\textwidth}{\centering \textit{[Figure~1 placeholder]}\\[4pt]
TorchCell overview --- 4 panels: (a) data-to-KG flow; (b) multimodal reference
cell; (c) perturbation operator; (d) CGT teaser. Replace with exported vector PDF.}}
\caption{TorchCell: ontology-unified knowledge graph and the perturbation-operator
cell representation. Panels (a)--(d) per outline.}
\label{fig:overview}
\end{figure*}

%% ===== Discussion (~500-650 words) =====
\section{Discussion}\label{sec:discussion}
\textit{[FILLER -- Discussion.]} A strain-aware, biologically grounded
representation is the differentiator; limitations include single-organism scope,
morphology-data maturity, and error-bar provenance.

%% ===== Methods (no word limit; full math here) =====
\section{Methods}\label{sec:methods}
\textit{[FILLER -- Methods.]} Ontology and knowledge-graph construction; dataset
and perturbation-operator formalism; CGT architecture and the graph-regularized
attention loss; training, splits, metrics; baselines (DANGO, DCell, Yeast9 FBA);
inference at scale; statistics and error-bar definitions.

\backmatter

\bmhead{Supplementary information}
% Supplementary Figures S1-S6 (see outline). S1 = classical-ML baselines.

\bmhead{Acknowledgments}

\section*{Declarations}
% Funding; competing interests; data availability (GitHub, NCSA Neo4j, HuggingFace);
% code availability; author contributions.

%% ===== Bibliography =====
\nocite{*}% FILLER: force all seeded refs to render in the preview; remove for submission
\bibliography{references}
